\scanpage{21}
\mspage{17}
% archive sheet number: 20

à partir de $\operatorname{Omb}^{\circ}(\Phi)=\Sigma$ par
\[
\left\{\begin{aligned}
\Phi_1&=\text{ensemble des éléments de }\deleted{\operatorname{Omb}^{\circ}(\Phi)}\Sigma\ \deleted{\unsure{tels}},\\
&\qquad\text{maximaux pour }\overset{\circ}{\ll},\text{ qui sont dans }\unsure{M_1|\mathcal M_1},\\
&\deleted{\Phi_0=\{s\in\unsure{\Sigma|M_0}\mid}\\
\Phi_{\mathrm{is}}&=\text{ensemble des éléments de }\Sigma,\text{ maximaux}\\
&\qquad\text{pour }\overset{\circ}{\ll},\text{ qui sont dans }M_0.\\
\Phi_0\setminus\Phi_{\mathrm{is}}&=\left(\bigcup_{X\in\Phi_1}\partial X\right)\cap\Sigma
 =\deleted{\{s\in\Sigma\cap M_0\}},\\
&=\{s\in\Sigma\cap M_0\mid\exists\,X\in\Phi_1\text{ t.q. }s\unsure{\triangleleft|\in\partial}X\}
\end{aligned}\right.
\]

On trouve que le cosupport d'une partie constructible de $M$ est encore constructible, et le passage au cosupport, \inserted{correspond} du côté des polygones, au passage à \deleted{\unsure{support}} le polygone complémentaire. Le passage au \inserted{support correspond \unsure{plutôt} au passage à l'\uline{enveloppe} (\uline{\unsure{irredondante|irrédundante}}) de la \uline{figure envisagée}.}
\margin{$\unsure{\flat|b}\Phi^{**}$}

Considérons un drapeau $T\in\operatorname{Drap}(I)$, $T\supset\partial I$,
\[
t_0=\omega_0<t_1\ldots<t_n=\omega_1
\]
($n=\text{longueur du drapeau}$), les strates constructibles sont
\[
\underset{\substack{\Vert\\\omega_0}}{\{t_0\}},\ \deleted{\unsure{\partial}}I_{t_0t_1},\ t_1,\ I_{t_1,t_2}\quad\cdots\quad I_{t_{\unsure{n-1}},t_n},\ \underset{\substack{\Vert\\\omega_1}}{t_n}.
\]
On appelle polygone \uline{subordonné} au drapeau $T$, toute partie \inserted{$\Phi$} de cet ensemble $\Phi_T$. On trouve alors (sans se préoccuper du support) ceci, $\forall\,\Phi\subset\Phi_T$, on désigne par $C\Phi$ son complémentaire dans $\Phi_T$
